\subsection{Gorenstein rings and flat algebras}

PROPOSITION 12. Let \(\rho: A \to B\) a local homomorphism of Noetherian local rings, making B a flat A-module. The following conditions are equivalent:

\begin{enumerate}
\def\labelenumi{(\roman{enumi})}
\item
  B is a Gorenstein ring;
\item
  A and \(\kappa_{\Lambda} \otimes_{\Lambda} B\) are Gorenstein rings.
\end{enumerate}

Let us first treat the case where the rings A and B are Artinian. Denote by C the local ring. \(\kappa_{\Lambda} \otimes_{\mathrm{A}} \mathrm{B}\) ; its residue field \(\kappa_{\mathrm{C}}\) is identified with \(\kappa_{\mathrm{B}}\) Since B is flat over A, the B-module \(\operatorname{Hom}_{\mathrm{B}}(\mathrm{C}, \mathrm{B})\) is isomorphic to \(\operatorname{Hom}_{\mathrm{A}}(\kappa_{\mathrm{A}}, \mathrm{A}) \otimes_{\mathrm{A}} \mathrm{B}\) (I, § 2, no. 10, prop. 11), hence to \(\operatorname{Hom}_{\mathrm{A}}(\kappa_{\mathrm{A}}, \mathrm{A}) \otimes_{\kappa_{\Lambda}} \mathrm{C}\) From this, we deduce a sequence of isomorphisms.

\[
\begin{array}{c} \operatorname{Hom} _ {\mathrm{B}} (\kappa_ {\mathrm{B}}, \mathrm{B}) \longrightarrow \operatorname{Hom} _ {\mathrm{C}} (\kappa_ {\mathrm{C}}, \operatorname{Hom} _ {\mathrm{B}} (\mathrm{C}, \mathrm{B})) \dashrightarrow \operatorname{Hom} _ {\mathrm{C}} (\kappa_ {\mathrm{C}}, \operatorname{Hom} _ {\Lambda} (\kappa_ {\mathrm{A}}, \mathrm{A}) \otimes_ {\kappa_ {\mathrm{A}}} \mathrm{C}) \\ \longrightarrow \operatorname{Hom} _ {\Lambda} (\kappa_ {\Lambda}, \mathrm{A}) \otimes_ {\kappa_ {\mathrm{A}}} \operatorname{Hom} _ {\mathrm{C}} (\kappa_ {\mathrm{C}}, \mathrm{C}). \end{array}
\]

In particular, we have \(\left[\mathrm{Hom}_{\mathrm{B}}(\kappa_{\mathrm{B}},\mathrm{B}):\kappa_{\mathrm{B}}\right] = \left[\mathrm{Hom}_{\mathrm{A}}(\kappa_{\mathrm{A}},\mathrm{A}):\kappa_{\mathrm{A}}\right]\left[\mathrm{Hom}_{\mathrm{C}}(\kappa_{\mathrm{C}},\mathrm{C}):\kappa_{\mathrm{C}}\right]\) and the proposition then follows from Lemma 1 of No. 7.

Let us turn to the general case. If in the statement one replaces the word “Gorenstein” by the word “Macaulay,” the proposition is a special case of Prop. 10 of § 2, no. 7. We may therefore assume that the rings A, B, and C = \(\kappa_{\mathrm{A}} \otimes_{\mathrm{A}} \mathrm{B}\) are Macaulay. The B-module C is Macaulay (§ 2, no. 1, example 4). Set \(r = \dim(\mathrm{A})\) , \(s = \dim(\mathrm{C})\) . There exists an A-regular sequence ( \(x_1, \ldots, x_r\) ) of elements of \(\mathfrak{m}_{\mathrm{A}}\) and a sequence ( \(y_1, \ldots, y_s\) ) of elements of \(\mathfrak{m}_{\mathrm{B}}\) regular for the B-module C (§ 2, no. 3, prop. 3); let \(x\) denote the ideal of A and \(\mathfrak{y}\) the ideal of B that they generate respectively. The sequence ( \(y_1, \ldots, y_s, \rho(x_1), \ldots, \rho(x_r)\) ) is B-regular (§ 1, no. 6, prop. 11) and the A-module B/ \(\mathfrak{y}\) is flat (loc. cit., prop. 10). The homomorphism of A/ \(x\) in B/(xB + \(\mathfrak{y}\) ) deduced from \(\rho\) Taking quotients therefore makes B/(xB + \(\mathfrak{y}\) ) a flat (A/x)-module and the ring \(\kappa_{\mathrm{A}/x} \otimes_{\mathrm{A}/x} \mathrm{B}/(xB + \mathfrak{y})\) is isomorphic to C/ \(\mathfrak{y}\) C. The proposition thus follows from the first part of the proof, taking into account example 3 of no. 7.

COROLLARY 1. -- Let \(\mathfrak{p}: A \to B\) a homomorphism of Noetherian rings making B a flat A-module. The following conditions are equivalent:

\begin{enumerate}
\def\labelenumi{(\roman{enumi})}
\item
  B is a Gorenstein ring;
\item
  (resp. (iii)) for every prime (resp. maximal) ideal \(\mathfrak{q}\) of B, the rings \(\mathrm{A}_{\mathfrak{p}^{-1}(\mathfrak{q})}\) and \(\kappa (\mathfrak{p}^{-1}(\mathfrak{q}))\otimes_{\mathrm{A}}\mathrm{B}\) are Gorenstein.
\end{enumerate}

If moreover B is a \(\Lambda\) -module that is faithfully flat, these conditions imply that A is a Gorenstein ring.

Let q be a prime ideal of B; set \(\mathfrak{p} = \rho^{-1}(\mathfrak{q})\) . The ring \(B_{q}\) , isomorphic to a ring of fractions of \(B_{p}\) is flat over \(A_{p}\) For \(B_{q}\) Let A be a Gorenstein ring; it is necessary and sufficient that the rings \(A_{p}\) and \(\kappa(\mathfrak{p}) \otimes_{A_{\mathfrak{p}}} B_{\mathfrak{q}}\) Let them be so (prop. 12).

\((\mathrm{i}) \Rightarrow (\mathrm{ii})\): Let \(\mathfrak{q}\) be a prime ideal of B and set \(\mathfrak{p} = \rho^{-1}(\mathfrak{q})\). If B is a Gorenstein ring, then so is \(B_{\mathfrak{q}}\); hence \(A_{\mathfrak{p}}\) and \(\kappa(\mathfrak{p}) \otimes_{A_{\mathfrak{p}}} B_{\mathfrak{q}}\) are Gorenstein by what precedes. By the remark in § 2, No. 6, every localization of \(\kappa(\mathfrak{p}) \otimes_{\mathrm{A}} B\) at a prime ideal is then a Gorenstein ring, which implies that \(\kappa(\mathfrak{p}) \otimes_{\mathrm{A}} B\) is a Gorenstein ring, whence (ii).

\begin{enumerate}
\def\labelenumi{(\roman{enumi})}
\setcounter{enumi}{1}
\tightlist
\item
  \(\Rightarrow\) (iii): this is clear.
\end{enumerate}

(iii)⇒(i): for every maximal ideal n of B, it follows from the beginning of the proof (applied with q = n) that \(B_{n}\) is a Gorenstein ring, whence (i).

If B is a faithfully flat A-module, the map \(^{q}\mathfrak{p}:\operatorname{Spec}(B)\longrightarrow\operatorname{Spec}(\Lambda)\) is surjective (II, § 2, n° 5, cor. 4 of prop. 11), whence the last assertion.

COROLLARY 2. Let A be a Noetherian ring, J an ideal of A, and  the separated completion of A for the J-adic topology. Consider the following conditions:

\begin{enumerate}
\def\labelenumi{(\roman{enumi})}
\item
  A is a Gorenstein ring;
\item
  \(\widehat{\mathrm{A}}\) is a Gorenstein ring;
\item
  for every maximal ideal m of A containing J, \(A_{m}\) is a Gorenstein ring ;
\item
  for every prime ideal p of A such that \(p + J \neq A\) , \(A_{p}\) and \(\kappa(\mathfrak{p}) \otimes_{A} \widehat{A}\) are Gorenstein rings.
\end{enumerate}

Conditions (ii) through (iv) are equivalent, and are implied by (i). When the ideal J is contained in the radical of A, conditions (i) through (iv) are equivalent.

We deduce this corollary from cor. 1 in the same way that we deduced corollary 4 from prop. 9 of § 2, no. 7: it suffices in the proof to replace the word “Macaulay” by the word “Gorenstein.”

COROLLARY 3.--- Let A be a Gorenstein ring and \((\mathrm{T}_{i})_{i\in\mathrm{I}}\) a finite family of indeterminates. The rings \(\mathrm{A}[(\mathrm{T}_{i})_{i\in\mathrm{I}}]\) and \(\mathrm{A}[((\mathrm{T}_{i})_{i\in\mathrm{I}}]]\) are Gorenstein rings.

For every field \(k\) , the ring \(k[\mathrm{T}]\) is Gorenstein (no. 7, example 5); moreover, the ring A{[}T{]} is Noetherian. It is therefore a Gorenstein ring (cor. 1). By induction on \(\operatorname{Card}(\mathrm{I})\) that \(\mathrm{A}[(\mathrm{T}_i)_{i\in \mathbf{I}}]\) is a Gorenstein ring; then we apply Cor. 2.

